\documentclass[master.tex]{subfiles}

\begin{document}

% ======================== TITLE ================================
\ifSubfilesClassLoaded{%
  % Standalone title block
  \begin{center}
    {\LARGE\bfseries\color{isublue} Chapter 2: DAGs and d-Separation}\\[6pt]
    {\large Theory and Methods in Causal Inference}\\[4pt]
    {\normalsize Department of Statistics, Iowa State University}
  \end{center}
  \bigskip
  \hrule height 1.2pt \relax
  \smallskip
}{%
  % Master book: chapter heading
  \chapter{DAGs and d-Separation}
}
% ======================== LEARNING OBJECTIVES ==================
\begin{tcolorbox}[defstyle, title={Learning Objectives}]
By the end of this chapter, students should be able to:
\begin{enumerate}
  \item Construct a DAG from a verbal description of a causal model and identify
        parents, descendants, ancestors, and non-descendants of any node.
  \item Classify every intermediate node on a path as a chain, fork, or collider,
        and apply the blocking rule for each.
  \item Apply the d-separation criterion to determine all conditional independence
        relationships implied by a DAG.
  \item Recognize and avoid collider bias, including the case where a descendant
        of a collider is conditioned upon.
  \item State the Markov factorization and use it to write the joint density of a
        DAG in terms of local conditional densities.
  \item Interpret d-separation results as conditional independence assumptions
        on observed variables, and recognize their role as the graphical
        expression of causal assumptions.
\end{enumerate}
\end{tcolorbox}

% ---------------------------------------------------------------
% Styles and macros specific to this chapter (not yet in master).
% Placed here (after the LO box, before §1) so the definition block
% does not inject blank-line spacing between \chapter{} and the box.
% The only \begin{remark} in this chapter is in §4, well below.
% (remarkstyle and remark environment are defined in causal_inference_master.tex)

% ==============================================================
\section{Directed Acyclic Graphs}
% ==============================================================

DAGs allow us to translate qualitative causal assumptions into quantitative
statistical restrictions.  This single idea underlies everything in this
chapter: once we draw a graph encoding what we believe about the causal
structure of a problem, the graph tells us exactly which variables are
independent, which are confounded, and what we need to condition on to
identify a causal effect.

\subsection{Basic Definition}

\begin{defn}{Directed Acyclic Graph}
A \emph{directed acyclic graph} (DAG) $\Gcal = (V, E)$ consists of a finite
set of nodes $V$ and a set of directed edges $E \subseteq V \times V$ such that there
is no directed cycle: no node $V_i$ has a directed path to itself.
\end{defn}

In a causal DAG each node represents a variable and each directed edge $A \to B$
encodes ``$A$ is a direct cause of $B$, given all other variables in the model.''
The acyclicity condition is not merely a technical convenience: it is necessary for
the Markov factorization $p(v_1,\dots,v_k) = \prod_i p(v_i \mid \Pa(v_i))$ to be
well-defined.  If the graph contained a cycle, say $A \to B \to A$, then
$p(a \mid \Pa(A)) = p(a \mid b)$ and $p(b \mid \Pa(B)) = p(b \mid a)$
would be mutually dependent with no consistent joint distribution to anchor them.

\subsection{Structural Relationships}

\begin{defn}{Structural relationships in a DAG}
For a node $v \in \mathcal{V}$ in $\Gcal = (\mathcal{V}, E)$:
\begin{itemize}
  \item $\Pa(v) = \{w \in \mathcal{V} : w \to v \in E\}$ are the \emph{parents} of $v$.
  \item $\De(v)$ = all nodes reachable from $v$ by directed paths
        are the \emph{descendants} of $v$.
  \item $\An(v)$ = all nodes with a directed path to $v$ are the
        \emph{ancestors} of $v$.
  \item $\Nd(v) = \mathcal{V} \setminus (\De(v) \cup \{v\})$ are the
        \emph{non-descendants} of $v$.
\end{itemize}
\end{defn}

\begin{example}{Education, Earnings, and Family Background}
Consider the causal story: education ($E$) affects earnings ($Y$); family
background ($B$) is a common cause of both; and neighborhood ($N$) affects
education but has no direct effect on earnings.

\begin{center}
\begin{tikzpicture}[node distance=1.8cm]
  \node[node] (N) {$N$};
  \node[node, right of=N] (E) {$E$};
  \node[node, right of=E] (Y) {$Y$};
  \node[node, below=1.0cm of E] (B) {$B$};
  \draw[edge] (N) -- (E);
  \draw[edge] (E) -- (Y);
  \draw[edge] (B) -- (E);
  \draw[edge] (B) -- (Y);
\end{tikzpicture}
\end{center}

The Markov factorization is $p(n,b,e,y) = p(n)\,p(b)\,p(e\mid n,b)\,p(y\mid e,b)$.
Reading off structural relationships: $\Pa(E) = \{N,B\}$,
$\Pa(Y) = \{E,B\}$, $\Pa(N)=\Pa(B)=\varnothing$.
The descendants of $N$ are $\{E,Y\}$; $B$ and $N$ are non-descendants of each other.
\end{example}

\begin{example}{The IV DAG}
The canonical instrumental variable model involves four nodes: $Z$ (instrument),
$T$ (treatment), $Y$ (outcome), and $U$ (unobserved confounder).

\begin{center}
\begin{tikzpicture}[node distance=1.8cm]
  \node[node] (Z) {$Z$};
  \node[node, right of=Z] (T) {$T$};
  \node[node, right of=T] (Y) {$Y$};
  \node[unode, below=1.0cm of T] (U) {$U$};
  \draw[edge] (Z) -- (T);
  \draw[edge] (T) -- (Y);
  \draw[dedge] (U) -- (T);
  \draw[dedge] (U) -- (Y);
\end{tikzpicture}
\end{center}

$\Pa(T) = \{Z,U\}$, $\Pa(Y) = \{T,U\}$,
$\Pa(Z) = \Pa(U) = \varnothing$.
The descendants of $Z$ are $\{T, Y\}$; $U$ and $Z$ are non-descendants of each other.
\end{example}

\begin{example}{The Mediation DAG}
In a mediation model, treatment $T$ affects outcome $Y$ both directly and
through an intermediate variable $M$ (the mediator).  An unobserved confounder
$U$ creates a back-door path $T \leftarrow U \rightarrow Y$, making $T$
endogenous.  The mediator $M$ is itself unconfounded.

\begin{center}
\begin{tikzpicture}[node distance=1.8cm]
  \node[node] (T) {$T$};
  \node[node, right of=T] (M) {$M$};
  \node[node, right of=M] (Y) {$Y$};
  \node[unode, below=1.0cm of T] (U) {$U$};
  \draw[edge] (T) -- (M);
  \draw[edge] (M) -- (Y);
  \draw[edge] (T) to[bend left=30] (Y);
  \draw[dedge] (U) -- (T);
  \draw[dedge] (U) -- (Y);
\end{tikzpicture}
\end{center}

$\Pa(T) = \{U\}$, $\Pa(M) = \{T\}$,
$\Pa(Y) = \{M, T, U\}$, $\Pa(U) = \varnothing$.
The total effect of $T$ on $Y$ travels along two paths: the \emph{direct}
path $T \to Y$ and the \emph{indirect} path $T \to M \to Y$.
Because $M$ is unconfounded, this structure satisfies the front-door criterion
(Chapter~3): the causal effect of $T$ on $Y$ is identifiable despite the
unmeasured confounder $U$.
\end{example}

% ==============================================================
\section{Paths, Blocking, and d-Separation}
% ==============================================================

A DAG encodes direct causal relationships through its edges, but variables
can also be statistically related through longer routes that mix directed
and reversed edges.  The concept of a \emph{path} formalises these routes,
and the question of whether a path transmits or blocks statistical
dependence --- depending on which variables are conditioned upon --- is
precisely what d-separation answers.  The key insight is that conditioning
can both \emph{close} channels that were open (by blocking a chain or fork)
and \emph{open} channels that were closed (by activating a collider).
Figure~\ref{w2:fig:chain-fork-collider} illustrates all three cases.

\subsection{Path Structures}

\begin{defn}{Path}
A \emph{path} between nodes $X$ and $Y$ in $\Gcal$ is any sequence of
distinct nodes $X = V_0, V_1, \dots, V_k = Y$ such that each consecutive pair
is connected by an edge in either direction.
\end{defn}

Every intermediate node $V_m$ on a path plays one of three structural roles.
In a \textbf{chain} ($V_{m-1} \to V_m \to V_{m+1}$), $V_m$ is a causal
intermediary through which influence is transmitted.  In a \textbf{fork}
($V_{m-1} \leftarrow V_m \to V_{m+1}$), $V_m$ is a common cause that
creates confounding between the two endpoints.  In a \textbf{collider}
($V_{m-1} \to V_m \leftarrow V_{m+1}$), both arrows point \emph{into}
$V_m$; this is the crucial asymmetric case whose behavior under
conditioning is the opposite of the other two, as
Figure~\ref{w2:fig:chain-fork-collider} shows.

\begin{figure}[t]
\centering
\resizebox{\textwidth}{!}{%
\begin{tikzpicture}[>=stealth, node distance=1.6cm and 1.8cm]
  %% -------- Chain --------
  \node[node] (A1) at (0,0)   {$A$};
  \node[node] (M1) at (1.8,0) {$M$};
  \node[node] (B1) at (3.6,0) {$B$};
  \draw[edge] (A1) -- (M1);
  \draw[edge] (M1) -- (B1);
  \node[font=\small\bfseries\color{isublue}] at (1.8, 0.9) {Chain};
  \node[font=\small\color{darkgrey}] at (1.8,-0.9)
    {$A \nindep B$
     \quad but \quad
     $A \indep B \mid M$};
  %% -------- Fork --------
  \node[node] (A2) at (6.2,0) {$A$};
  \node[node] (M2) at (8.0,0) {$M$};
  \node[node] (B2) at (9.8,0) {$B$};
  \draw[edge] (M2) -- (A2);
  \draw[edge] (M2) -- (B2);
  \node[font=\small\bfseries\color{isublue}] at (8.0, 0.9) {Fork};
  \node[font=\small\color{darkgrey}] at (8.0,-0.9)
    {$A \nindep B$
     \quad but \quad
     $A \indep B \mid M$};
  %% -------- Collider --------
  \node[node] (A3) at (12.4,0) {$A$};
  \node[node] (M3) at (14.2,0) {$M$};
  \node[node] (B3) at (16.0,0) {$B$};
  \draw[edge] (A3) -- (M3);
  \draw[edge] (B3) -- (M3);
  \node[font=\small\bfseries\color{isublue}] at (14.2, 0.9) {Collider};
  \node[font=\small\color{darkgrey}] at (14.2,-0.9)
    {$A \indep B$
     \quad but \quad
     $A \nindep B \mid M$};
  %% -------- Vertical dividers --------
  \draw[dotted, color=accent!40, thin] (5.0, 1.4) -- (5.0,-1.4);
  \draw[dotted, color=accent!40, thin] (11.0, 1.4) -- (11.0,-1.4);
\end{tikzpicture}%
}
\caption{The three fundamental path structures in a DAG, with their default
independence status and the effect of conditioning on the middle node $M$.
Chains and forks are \emph{open by default} and are blocked by conditioning
on $M$.  Colliders are \emph{closed by default} and are opened by
conditioning on $M$ (or any descendant of $M$).}
\label{w2:fig:chain-fork-collider}
\end{figure}

\subsection{Blocking Rules}

Conditioning on a node affects the paths it lies on differently depending on
its structural role:

\begin{center}
\renewcommand{\arraystretch}{1.35}
\begin{tabular}{@{}lllp{5.2cm}@{}}
\toprule
\textbf{Structure} & \textbf{Pattern} & \textbf{Effect of conditioning} & \textbf{Key intuition} \\
\midrule
Chain    & $A \to M \to B$    & Blocks the path & Causal transmission \\
Fork     & $A \leftarrow M \to B$ & Blocks the path & Common cause \\
Collider & $A \to M \leftarrow B$ & \textbf{Opens} the path & Common effect \\
\bottomrule
\end{tabular}
\end{center}

\medskip
The critical asymmetry --- colliders are \emph{closed by default} and
\emph{opened} by conditioning, while chains and forks are \emph{open by
default} and \emph{closed} by conditioning --- is what makes collider bias
so easy to overlook in practice.  The following three examples anchor each
rule to a concrete story.

\medskip
\noindent\textbf{Chain: Smoking $\to$ Tar $\to$ Cancer.}
Let $A$ = smoking, $M$ = tar deposits in the lungs, $B$ = lung cancer.
The path $A \to M \to B$ is a chain: smoking causes tar accumulation, which
in turn causes cancer.
Marginally, heavy smokers have elevated cancer rates --- the path is open.
But if we condition on tar level (e.g., restrict to patients with the same
measured tar burden), smoking no longer predicts cancer within that stratum:
once the tar level is fixed, the causal channel is fully accounted for.
Conditioning on the intermediary $M$ \emph{blocks} the path.

\medskip
\noindent\textbf{Fork: Poverty $\to$ Poor Diet; Poverty $\to$ Lack of Exercise.}
Let $M$ = poverty, $A$ = poor diet, $B$ = lack of exercise.
The path $A \leftarrow M \to B$ is a fork: poverty is a common cause of
both outcomes.
Unconditionally, diet quality and exercise are correlated --- not
because one causes the other, but because both are driven by poverty.
Conditioning on income level removes this spurious association: among people
at the same income, the diet--exercise correlation disappears.
Conditioning on the common cause $M$ \emph{blocks} the path.

\medskip
\noindent\textbf{Collider: Accident $\to$ hospitalization $\leftarrow$ Cancer.}
Let $A$ = traffic accident, $B$ = cancer, $M$ = hospitalization.
Both accidents and cancer independently cause hospitalization, so the path
$A \to M \leftarrow B$ is a collider at $M$.
In the general population, having a traffic accident tells you nothing about
your cancer risk: $A \indep B$.
But among hospitalised patients --- conditioning on $M$ --- the two become
negatively associated.
If we learn that a hospitalised patient was \emph{not} in a traffic accident,
cancer or another serious illness becomes more likely as an explanation for
the admission.
Conditioning on the common effect $M$ \emph{opens} the path.
This is Berkson's bias \citep{berkson1946limitations}.

\subsection{The d-Separation Criterion}

\begin{defn}{d-Blocking and d-Separation}
\label{w2:def:dsep}
A path $\pi$ is \emph{d-blocked} by a set $\mathbf{S}$ if:
\begin{enumerate}
  \item $\pi$ contains a chain $A \to M \to B$ or fork $A \leftarrow M \to B$
        with $M \in \mathbf{S}$; or
  \item $\pi$ contains a collider $A \to C \leftarrow B$ with $C \notin
        \mathbf{S}$ and no descendant of $C$ is in $\mathbf{S}$.
\end{enumerate}
Nodes $X$ and $Y$ are \emph{d-separated} by $\mathbf{S}$, written
$(X \indep Y \mid \mathbf{S})_{\Gcal}$, if every path between
$X$ and $Y$ in $\Gcal$ is d-blocked by $\mathbf{S}$.
\end{defn}

\begin{example}{Applying Definition~\ref{w2:def:dsep} to the Three Toy Settings}
\label{w2:ex:threesettings}
We revisit the three examples from Section~2.2 and verify each blocking
claim directly using Definition~\ref{w2:def:dsep}.

\medskip
\noindent\textbf{(i) Chain --- Smoking $\to$ Tar $\to$ Cancer.}
Let $A$ = smoking, $M$ = tar, $B$ = cancer.
The only path between $A$ and $B$ is the chain $A \to M \to B$.

\begin{itemize}
  \item $\mathbf{S} = \varnothing$: the path contains no conditioned-on
        non-collider, so clause~(1) does not apply.
        The collider clause~(2) is also irrelevant (there is no collider).
        The path is \emph{not} d-blocked.
        $\Rightarrow A \nindep B$ (smokers have higher cancer rates).
  \item $\mathbf{S} = \{M\}$: the chain node $M \in \mathbf{S}$,
        so clause~(1) applies and the path is d-blocked.
        $\Rightarrow A \indep B \mid M$ (at fixed tar level,
        smoking no longer predicts cancer).
\end{itemize}

\medskip
\noindent\textbf{(ii) Fork --- Poverty $\to$ Poor Diet; Poverty $\to$ Lack of Exercise.}
Let $M$ = poverty, $A$ = poor diet, $B$ = lack of exercise.
The only path is the fork $A \leftarrow M \to B$.

\begin{itemize}
  \item $\mathbf{S} = \varnothing$: no non-collider on the path is in
        $\mathbf{S}$, so the path is not d-blocked.
        $\Rightarrow A \nindep B$ (poor diet and lack of
        exercise are correlated through shared poverty).
  \item $\mathbf{S} = \{M\}$: the fork node $M \in \mathbf{S}$, so
        clause~(1) applies and the path is d-blocked.
        $\Rightarrow A \indep B \mid M$ (at fixed income level,
        the diet--exercise association disappears).
\end{itemize}

\medskip
\noindent\textbf{(iii) Collider --- Accident $\to$ hospitalization $\leftarrow$ Cancer.}
Let $A$ = accident, $M$ = hospitalization, $B$ = cancer.
The only path is the collider $A \to M \leftarrow B$.

\begin{itemize}
  \item $\mathbf{S} = \varnothing$: the path contains collider $M$
        with $M \notin \mathbf{S}$ and no descendant of $M$ in
        $\mathbf{S}$, so clause~(2) applies and the path is d-blocked.
        $\Rightarrow A \indep B$ (accidents and cancer are
        independent in the general population).
  \item $\mathbf{S} = \{M\}$: now $M \in \mathbf{S}$, so clause~(2) no
        longer blocks the path.  No other clause applies, so the path is
        \emph{open}.
        $\Rightarrow A \nindep B \mid M$ (among hospitalised
        patients, ruling out an accident raises the probability of cancer ---
        Berkson's bias).
\end{itemize}
\end{example}

Example~\ref{w2:ex:threesettings} showed that d-separation is a purely graphical
calculation.  The following theorem is what gives it statistical force: it
guarantees that every d-separation statement in the graph corresponds to a
genuine conditional independence in the joint distribution.

\begin{thm}{Verma \& Pearl, 1988; \citep{pearl2009causality}, Theorem~1.2.4}
\label{w2:thm:dsep}
If $(X \indep Y \mid \mathbf{S})_{\Gcal}$, then
$X \indep Y \mid \mathbf{S}$ in every distribution that is Markov
with respect to $\Gcal$.
\end{thm}

In other words, if we can show from the graph alone that every path
between $X$ and $Y$ is d-blocked by $\mathbf{S}$, we are entitled to
conclude that $X$ and $Y$ are conditionally independent in the data ---
without any calculation involving the actual distribution.
This is the theorem's practical value: it turns a visual inspection of
the graph into a licence to drop terms from the joint density, to omit
variables from a regression, or to assert that a treatment assignment is
unconfounded given a set of covariates.
Example~\ref{w2:ex:threesettings}(iii) illustrates the contrapositive: when the
path is \emph{open} (conditioning on the collider $M$), the theorem gives
no independence guarantee, and indeed the dependence
$A \nindep B \mid M$ materialises in the data as Berkson's bias.

Table~\ref{w2:tab:markov-implications} summarizes the Markov factorization and
its independence implication for each of the three path structures.

\begin{center}
\renewcommand{\arraystretch}{1.6}
\begin{tabular}{@{}lll@{}}
\toprule
\textbf{Structure} & \textbf{Markov factorization} $p(a,m,b)$
  & \textbf{d-separation} \\
\midrule
Chain $A\!\to\!M\!\to\!B$
  & $p(a)\,p(m\mid a)\,p(b\mid m)$
  & $(A\indep B\mid M)_{\Gcal}$ \\[4pt]
Fork $A\!\leftarrow\!M\!\to\!B$
  & $p(m)\,p(a\mid m)\,p(b\mid m)$
  & $(A\indep B\mid M)_{\Gcal}$ \\[4pt]
Collider $A\!\to\!M\!\leftarrow\!B$
  & $p(a)\,p(b)\,p(m\mid a,b)$
  & $(A\indep B)_{\Gcal}$ \\
\bottomrule
\end{tabular}
\captionof{table}{Markov factorization and d-separation result for each path
structure.  In the chain and fork, the factor $p(b\mid m)$ depends on
$b$ only through $m$, making $A\indep B\mid M$ immediate.
In the collider, the factors $p(a)$ and $p(b)$ are free of each other,
so marginalising over $M$ yields $p(a,b)=p(a)\,p(b)$.}
\label{w2:tab:markov-implications}
\end{center}

\begin{example}{Theorem~\ref{w2:thm:dsep} via the Markov factorization --- the Chain}
We verify Theorem~\ref{w2:thm:dsep} for the chain $A \to M \to B$ with
$\mathbf{S} = \{M\}$.
The Markov factorization gives
\[
  p(a, m, b) \;=\; p(a)\,p(m \mid a)\,p(b \mid m).
\]
Conditioning on $M = m$:
\[
  p(a, b \mid m)
  \;=\; \frac{p(a, m, b)}{p(m)}
  \;=\; \frac{p(a)\,p(m \mid a)\,p(b \mid m)}{p(m)}
  \;=\; p(a \mid m)\,p(b \mid m).
\]
The joint conditional density factorises into a product of the two
marginal conditionals, which is exactly $A \indep B \mid M$ in
the distribution.  This confirms Theorem~\ref{w2:thm:dsep} for the chain: the
d-separation statement $(A \indep B \mid M)_{\Gcal}$
from Example~\ref{w2:ex:threesettings}(i) implies the distributional independence.

The fork $A \leftarrow M \to B$ with $\mathbf{S} = \{M\}$ follows by an
identical calculation using $p(a, m, b) = p(m)\,p(a \mid m)\,p(b \mid m)$.
The collider case requires a different argument and is left as Problem~5.
\end{example}

\subsection{Practical Ways to Check d-Separation}

The formal definition of d-separation can appear abstract when applied to
larger graphs.  Two complementary algorithmic viewpoints are useful in
practice, both equivalent to the definition.

\paragraph{1. Bayes-Ball Intuition.}
Imagine releasing a ball from $X$ and asking whether it can reach $Y$,
given that nodes in $\mathbf{S}$ are observed (shaded).
The ball follows four local rules:
\begin{itemize}
  \item At a chain or fork node \emph{not} in $\mathbf{S}$: ball passes
        through.
  \item At a chain or fork node in $\mathbf{S}$: ball stops.
  \item At a collider not in $\mathbf{S}$ (and no descendant in
        $\mathbf{S}$): ball stops.
  \item At a collider in $\mathbf{S}$ (or with a descendant in
        $\mathbf{S}$): ball passes through.
\end{itemize}
If the ball can reach $Y$ from $X$, then $X \nindep Y \mid
\mathbf{S}$.  If every path is stopped, then $X \indep Y \mid
\mathbf{S}$.  The algorithm is formalised in \citet{shachter1998bayesball}.

\paragraph{2. Moral Graph Transformation.}
There is also a graph-transformation approach, which connects DAG models
to undirected graphical models (Markov random fields).
\begin{enumerate}
  \item Restrict the DAG to the \emph{ancestral set}
        $\An(X \cup Y \cup \mathbf{S})$: discard all nodes that
        are not ancestors of $X$, $Y$, or any element of $\mathbf{S}$.
  \item \emph{Moralise} the graph: for every collider $A \to C \leftarrow
        B$ in the ancestral subgraph, add an undirected edge $A - B$
        (``marry the parents'').
  \item Remove all edge directions to obtain an undirected graph.
  \item Delete all nodes in $\mathbf{S}$ (and their incident edges).
\end{enumerate}
If $X$ and $Y$ are disconnected in the resulting graph, then
$(X \indep Y \mid \mathbf{S})_{\Gcal}$.
The moralization step is what makes the collider rule visible in the
undirected representation: without adding the $A$--$B$ edge, the path
through $C$ would appear disconnected even when $C$ is conditioned upon,
which would give the wrong answer.  A full treatment is in
\citet[][Ch.~3]{lauritzen1996graphical}.

\medskip
These two viewpoints --- Bayes-ball propagation and moral graph separation
--- are equivalent algorithmic ways of evaluating the d-separation criterion.
Students are encouraged to use whichever they find most natural, and to
cross-check with the other when in doubt.

\begin{example}{Checking d-Separation in a Four-Node DAG}
\label{w2:ex:fournode}
Consider the DAG with edges $Z \to T$, $U \to T$, $T \to Y$, $U \to Y$.

\begin{center}
\begin{tikzpicture}[node distance=1.8cm]
  \node[node] (Z) {$Z$};
  \node[node, right of=Z] (T) {$T$};
  \node[node, right of=T] (Y) {$Y$};
  \node[node, below=1.0cm of T] (U) {$U$};
  \draw[edge] (Z) -- (T);
  \draw[edge] (T) -- (Y);
  \draw[edge] (U) -- (T);
  \draw[edge] (U) -- (Y);
\end{tikzpicture}
\end{center}

We apply both methods to two queries.

\medskip
\noindent\textbf{Query 1.} Is $(Z \indep U)_{\Gcal}$,
i.e.\ $\mathbf{S} = \varnothing$?

\medskip
\noindent\emph{Bayes-ball.}
The only path from $Z$ to $U$ is $Z \to T \leftarrow U$.
Node $T$ is a collider with $T \notin \mathbf{S}$ and no descendant of $T$
in $\mathbf{S}$, so the ball is stopped at $T$.
No ball can reach $U$, hence $(Z \indep U)_{\Gcal}$.

\medskip
\noindent\emph{Moral graph.}
\begin{enumerate}
  \item \emph{Ancestral set} of $\{Z, U\} \cup \mathbf{S} = \{Z, U\}$:
        neither $Z$ nor $U$ has any parents, so
        $\An(\{Z,U\}) = \{Z, U\}$.
        Nodes $T$ and $Y$ are discarded.
  \item \emph{Moralise} the subgraph $\{Z, U\}$: it has no edges and no
        colliders, so nothing is added.
  \item \emph{Undirected graph}: isolated nodes $Z$ and $U$.
  \item \emph{Delete} $\mathbf{S} = \varnothing$: nothing changes.
\end{enumerate}
$Z$ and $U$ are disconnected $\Rightarrow$ $(Z \indep U)_{\Gcal}$.
Note that the key step is restricting to the ancestral set: this removes $T$
from the graph before moralization, so the collider $T$ never has a chance
to create a moral edge between $Z$ and $U$.

\medskip
\noindent\textbf{Query 2.} Is $(Z \indep U \mid T)_{\Gcal}$,
i.e.\ $\mathbf{S} = \{T\}$?

\medskip
\noindent\emph{Bayes-ball.}
Same path $Z \to T \leftarrow U$.
Now $T \in \mathbf{S}$: the collider is observed, so the ball passes through $T$.
The path is open, hence $Z \nindep U \mid T$.

\medskip
\noindent\emph{Moral graph.}
\begin{enumerate}
  \item \emph{Ancestral set} of $\{Z, U\} \cup \mathbf{S} = \{Z, U, T\}$:
        parents of $T$ are $Z$ and $U$, both already in the set.
        $\An(\{Z,U,T\}) = \{Z, U, T\}$; node $Y$ is discarded.
  \item \emph{Moralise}: the subgraph contains the collider $Z \to T
        \leftarrow U$, so add the moral edge $Z - U$.
  \item \emph{Undirected graph}: edges $Z - T$, $U - T$, $Z - U$.
  \item \emph{Delete} $\mathbf{S} = \{T\}$ (and its incident edges $Z-T$, $U-T$).
        Remaining graph: single edge $Z - U$.
\end{enumerate}
$Z$ and $U$ are connected $\Rightarrow$ $Z \nindep U \mid T$.
The moralization step is decisive: it adds the edge $Z - U$ \emph{before}
$T$ is deleted, ensuring the dependence induced by conditioning on the
collider is visible in the undirected graph.

\medskip
Comparing the two queries: $Z$ and $U$ are marginally independent (the
instrument is exogenous) but become dependent once we condition on $T$.
This is Berkson's bias, treated in full in Section~\ref{w2:sec:ivdag}.
\end{example}

% ==============================================================
\section{Collider Bias and the IV DAG}
\label{w2:sec:ivdag}
% ==============================================================

\subsection{Berkson's Bias}

\begin{warn}{Warning: Collider Bias \citep{berkson1946limitations}}
Conditioning on a collider variable induces a spurious association between
its parents, \emph{even when they are marginally independent}. The same effect
occurs when conditioning on any descendant of a collider.
\end{warn}

\begin{example}{Collider Bias in the IV DAG}
In the IV DAG, $T$ is a collider on the path $Z \to T \leftarrow U$.
Marginally, $Z \indep U$ (the instrument is exogenous). However,
\[
  Z \indep U \qquad \text{but} \qquad
  Z \nindep U \mid T
\]
(established graphically in Example~\ref{w2:ex:fournode}).
Conditioning on $T$ --- restricting to the treated subpopulation --- makes
the instrument $Z$ appear associated with $U$,
invalidating the IV argument within that stratum.
\end{example}

\subsection{Full d-Separation Analysis of the IV DAG}

We work through the IV DAG with edges $Z \to T$, $T \to Y$, $U \to T$,
$U \to Y$.  This is the same graph as in Example~\ref{w2:ex:fournode}, but now
$U$ is treated as \emph{unobserved}.  That assumption is precisely what earns
the name ``IV DAG'': because $U$ cannot be conditioned on, the back-door path
$T \leftarrow U \to Y$ cannot be blocked by adjustment, and the instrument
$Z$ becomes the only route to identification.

Before listing the paths, it is worth classifying them by \emph{type},
since students often conflate causal and associational dependence.

\begin{center}
\renewcommand{\arraystretch}{1.35}
\begin{tabular}{@{}llll@{}}
\toprule
\textbf{Endpoints} & \textbf{Path} & \textbf{Type} & \textbf{Why} \\
\midrule
$Z \leftrightarrow Y$ & $Z \to T \to Y$ & \emph{Causal} & Directed from $Z$ to $Y$; carries the IV signal \\
$Z \leftrightarrow Y$ & $Z \to T \leftarrow U \to Y$ & \emph{Associational} & Collider at $T$; activates back-door via $U$ \\
$Z \leftrightarrow U$ & $Z \to T \leftarrow U$ & \emph{Associational} & Collider at $T$; dormant unless $T$ conditioned on \\
\bottomrule
\end{tabular}
\end{center}

\medskip
\noindent The key distinction: a \emph{causal path} follows all edges in
their directed orientation from source to target.
An \emph{associational path} travels against at least one arrow, or passes
through a collider, and therefore does not represent a mechanism by which
$Z$ produces $Y$ --- it can nevertheless generate a statistical association
if the collider is conditioned upon.
With this in mind, the four questions become much easier to diagnose.

\medskip
\noindent\textbf{Question 1.} Is $Z \indep Y \mid \varnothing$?

\noindent\emph{Causal path} $Z \to T \to Y$: open (chain, nothing conditioned on).
\emph{Associational path} $Z \to T \leftarrow U \to Y$: collider at $T$,
$T \notin \varnothing$ --- blocked.
One open path suffices for dependence.
$\Rightarrow Z \nindep Y$.
The association is entirely \emph{causal} here; the back-door path is dormant.

\medskip
\noindent\textbf{Question 2.} Is $Z \indep Y \mid T$?

\begin{itemize}
  \item \emph{Causal path} $Z \to T \to Y$: chain with $T$ conditioned on
        --- \emph{blocked}. The causal channel is shut.
  \item \emph{Associational path} $Z \to T \leftarrow U \to Y$: collider
        at $T$ with $T \in \{T\}$ --- \emph{opened}. The back-door channel
        is activated.
\end{itemize}
$\Rightarrow Z \nindep Y \mid T$.
This is the key lesson: conditioning on $T$ simultaneously closes the causal
path \emph{and} opens the associational path, so the residual association
between $Z$ and $Y$ given $T$ is entirely spurious --- a collider artefact,
not a causal effect.

\medskip
\noindent\textbf{Question 3.} Is $Z \indep Y \mid \{T, U\}$?

\begin{itemize}
  \item \emph{Causal path} $Z \to T \to Y$: blocked by $T$.
  \item \emph{Associational path} $Z \to T \leftarrow U \to Y$: collider
        at $T$ opened (since $T \in \{T,U\}$), but the continuing segment
        $U \to Y$ is then blocked because $U \in \{T,U\}$.
\end{itemize}
$\Rightarrow Z \indep Y \mid \{T, U\}$.
Both the causal and the associational channel are closed; $Z$ carries no
information about $Y$ once treatment and the confounder are held fixed.

\medskip
\noindent\textbf{Question 4.} Is $Z \indep U \mid \varnothing$?

\noindent The only path, $Z \to T \leftarrow U$, is \emph{associational}
(collider at $T$) and is blocked since $T \notin \varnothing$.
$\Rightarrow Z \indep U$.
This is the \emph{exogeneity} assumption: the instrument is marginally
independent of the confounder precisely because the only path linking them
is a dormant collider path.

\medskip
Taken together, the four questions give a complete graphical account of the
three IV assumptions and one critical warning.
Question~1 establishes \emph{relevance}: $Z$ has an open causal path to $Y$
through $T$, so the instrument is not inert.
Question~4 establishes \emph{exogeneity}: the only path from $Z$ to the
confounder $U$ is a dormant collider path, so the instrument is
independent of unmeasured confounding.
Question~3 establishes the \emph{exclusion restriction}: once both $T$ and
$U$ are held fixed, $Z$ carries no residual information about $Y$, meaning
the instrument affects the outcome only through the treatment channel.
Note, however, that this conditioning requires observing $U$ --- which is
unavailable by assumption. The IV strategy exploits the exclusion restriction
indirectly, through the moment condition $\E[\varepsilon \cdot Z] = 0$ derived
in Chapter~7, without ever observing $U$.
Finally, Question~2 is the warning: conditioning on $T$ alone --- as a naive
analyst might do --- simultaneously closes the causal channel and opens the
confounded channel, producing a spurious association that is worse than no
adjustment at all.

% ==============================================================
\section{The Markov Property and factorization}
\label{w2:sec:markov}
% ==============================================================

\begin{defn}{Local Markov Property}
A distribution $P$ satisfies the \emph{local Markov property} with respect
to $\Gcal$ if, for every node $v_i \in \mathcal{V}$:
\[
  v_i \indep \Nd(v_i) \mid \Pa(v_i).
\]
That is, each node is conditionally independent of all its non-descendants
given its parents.
\end{defn}

\begin{remark}
The \emph{global Markov property} is the statement that every
d-separation in $\Gcal$ implies a conditional independence in $P$
--- precisely the content of Theorem~\ref{w2:thm:dsep}.
For DAGs, the local and global Markov properties are equivalent
\citep[][Proposition~3.27]{lauritzen1996graphical}: the local property
implies the global one via the Markov factorization, which is what is
actually used in identification arguments.
\end{remark}

An immediate consequence is the \emph{Markov factorization}:
\begin{equation}
  p(v_1, \dots, v_k) \;=\; \prod_{i=1}^{k} p\!\left(v_i \mid \Pa(v_i)\right).
  \label{w2:eq:markov}
\end{equation}
This factorization is the bridge between causal structure and probability:
knowing the parents of each node is sufficient to write the full joint density.

\begin{example}{Markov factorization in the Education Example}
For the DAG with edges $N \to E$, $B \to E$, $B \to Y$, $E \to Y$:
\[
  p(n, b, e, y) = p(n)\,p(b)\,p(e \mid n, b)\,p(y \mid e, b).
\]
A regression of $Y$ on $E$ alone is confounded by $B$. Conditioning on $B$
blocks the back-door path $E \leftarrow B \to Y$ and identifies the causal
effect $P(y \mid \doop(E{=}e))$ via the back-door formula (Chapter~3).
\end{example}

\begin{example}{Markov factorization in the IV DAG}
Including the latent $U$:
\[
  p(z, t, y, u) = p(z)\,p(u)\,p(t \mid z, u)\,p(y \mid t, u).
\]
Since $U$ is unobserved, the observable factorization is obtained by
marginalising:
\[
  p(z, t, y) = \int p(z)\,p(u)\,p(t \mid z, u)\,p(y \mid t, u)\,\mathrm{d}u.
\]
This cannot be simplified to $p(z)\,p(t\mid z)\,p(y\mid t)$ when $U$ is a
common cause of $T$ and $Y$: this is precisely the endogeneity problem.
\end{example}

% ==============================================================
\section{Worked Example: The Education--Earnings DAG}
\label{w2:sec:education}
% ==============================================================

This section traces all four steps of the pipeline --- graph structure,
Markov factorization, d-separation, and identification --- for the
Education--Earnings DAG introduced in Section~1.  It is intended as a
self-contained reference example that consolidates the tools of this chapter.

\medskip
\noindent\textbf{The causal story.}
Education ($E$) affects earnings ($Y$); family background ($B$) is a
common cause of both education and earnings; neighborhood ($N$) affects
education but has no direct effect on earnings.  All four variables are
observed; there are no latent confounders.

\begin{center}
\begin{tikzpicture}[node distance=1.8cm]
  \node[node] (N) {$N$};
  \node[node, right of=N] (E) {$E$};
  \node[node, right of=E] (Y) {$Y$};
  \node[node, below=1.0cm of E] (B) {$B$};
  \draw[edge] (N) -- (E);
  \draw[edge] (E) -- (Y);
  \draw[edge] (B) -- (E);
  \draw[edge] (B) -- (Y);
\end{tikzpicture}
\end{center}

\medskip
\noindent\textbf{Step 1 --- Markov factorization.}
The parents are $\Pa(N) = \Pa(B) = \varnothing$,
$\Pa(E) = \{N, B\}$, $\Pa(Y) = \{E, B\}$.
The Markov factorization is therefore
\[
  p(n, b, e, y) \;=\; p(n)\,p(b)\,p(e \mid n, b)\,p(y \mid e, b).
\]
Every factor on the right-hand side is a conditional density of an
observed variable given observed parents, so in principle each factor
is estimable from data.  This is the starting point for all
identification arguments: the full joint is expressed in terms of
quantities we can estimate.

\medskip
\noindent\textbf{Step 2 --- d-Separation.}
There are two paths between $N$ and $Y$:
\begin{itemize}
  \item Path 1: $N \to E \to Y$ (chain through $E$).
  \item Path 2: $N \to E \leftarrow B \to Y$ (collider at $E$, then fork leg $B\to Y$).
\end{itemize}
We examine four queries.

\smallskip
\emph{Query (i): Is $(N \indep Y)_{\Gcal}$?}
Path~1 is a chain with nothing conditioned on --- open.
Path~2 has collider $E$ with $E \notin \varnothing$ --- blocked.
One open path suffices: $N \nindep Y$.

\smallskip
\emph{Query (ii): Is $(N \indep Y \mid E)_{\Gcal}$?}
Path~1: chain with $E \in \{E\}$ --- blocked.
Path~2: collider $E$ with $E \in \{E\}$ --- \emph{opened}.
The activated path $N \to E \leftarrow B \to Y$ continues through $B$,
which is not in $\{E\}$, so the path is open.
Hence $N \nindep Y \mid E$.
This is a collider trap: conditioning on $E$ to ``control for education''
inadvertently opens a spurious association between $N$ and $Y$ through $B$.

\smallskip
\emph{Query (iii): Is $(N \indep Y \mid \{E, B\})_{\Gcal}$?}
Path~1: blocked by $E$.
Path~2: collider $E$ opened (since $E \in \{E,B\}$), but the continuing
segment $B \to Y$ has $B \in \{E,B\}$ --- blocked.
Both paths are blocked: $(N \indep Y \mid E, B)_{\Gcal}$.

\smallskip
\emph{Query (iv): Is $(N \indep B)_{\Gcal}$?}
The only path is $N \to E \leftarrow B$: collider at $E$ with
$E \notin \varnothing$ --- blocked.
Hence $N \indep B$.

\medskip
\noindent\textbf{Step 3 --- Conditional Independence.}
By Theorem~\ref{w2:thm:dsep}, the d-separation results translate into
distributional statements:
\[
  N \indep B, \qquad
  N \indep Y \mid E, B, \qquad
  N \nindep Y, \qquad
  N \nindep Y \mid E.
\]
The first two are the independence statements; the last two are the
dependences.  Three of the four are testable from data and could in
principle be used to validate or falsify the graph.
The statement $N \nindep Y \mid E$ is especially important
as a warning to practitioners: a researcher who conditions only on
education when studying the neighborhood--earnings association will
\emph{introduce} bias through the activated collider path, not remove it.

\medskip
\noindent\textbf{Step 4 --- Identification.}
We wish to identify $P(y \mid \doop(E{=}e))$, earnings under an
intervention that sets education to $e$.  The back-door criterion
(Chapter~3) requires a set $\mathbf{S}$ that (i) blocks every back-door
path from $E$ to $Y$ and (ii) contains no descendant of $E$.

The only back-door path is $E \leftarrow B \to Y$ (a fork at $B$).
Consider two candidate adjustment sets:
\begin{itemize}
  \item $\mathbf{S} = \{B\}$: blocks $E \leftarrow B \to Y$ (fork at
        $B$); $B$ is not a descendant of $E$. \emph{Valid.}
  \item $\mathbf{S} = \{N\}$: does not contain $B$, so the back-door
        path $E \leftarrow B \to Y$ remains open. \emph{Invalid.}
\end{itemize}
The second case illustrates why the criterion requires blocking
\emph{every} back-door path: $N$ is upstream of $E$ and appears
innocuous, but it does not close the confounding fork through $B$.

With $\mathbf{S} = \{B\}$, the back-door formula gives
\[
  P(y \mid \doop(E{=}e))
  \;=\; \sum_{b} P(y \mid e, b)\,P(b).
\]
Every term on the right involves only the observed distribution, so the
causal effect is identified entirely from observational data --- the
payoff of the graphical analysis.

% ==============================================================
\section{The Big Picture}
% ==============================================================

The concepts developed in this chapter form the graphical foundation for all
subsequent identification and estimation theory. The chain of reasoning runs
from a qualitative causal graph all the way to a formula for the
interventional distribution:

\bigskip
\begin{center}
\begin{tikzpicture}[
    box/.style = {
      rectangle, rounded corners=6pt,
      draw=isublue, fill=defbg, thick,
      text width=7.2cm, align=center,
      minimum height=1.1cm,
      font=\bfseries\color{isublue}
    },
    tool/.style = {
      rectangle, rounded corners=3pt,
      draw=accent, fill=white, thin,
      text width=4.6cm, align=center,
      font=\small\itshape\color{accent}
    },
    arr/.style  = {-{Stealth[length=7pt,width=5pt]}, line width=1.6pt, color=accent},
    node distance = 0.55cm
  ]
  % Four main boxes
  \node[box] (G)  {Graph Structure};
  \node[box, below=1.2cm of G]  (D)  {d-Separation};
  \node[box, below=1.2cm of D]  (CI) {Conditional Independence};
  \node[box, below=1.2cm of CI] (ID) {Identification Formula};

  % Arrows
  \draw[arr] (G)  -- (D);
  \draw[arr] (D)  -- (CI);
  \draw[arr] (CI) -- (ID);

  % Side annotations
  \node[tool, right=1.1cm of G]
    {DAG encodes direct causal relationships};
  \node[tool, right=1.1cm of D]
    {Read off all $\indep$ statements\\ from the graph};
  \node[tool, right=1.1cm of CI]
    {Test or exploit independence\\ structure in data};
  \node[tool, right=1.1cm of ID]
    {Back-door / front-door /\\ do-calculus (Chapter~3)};

  % Horizontal connectors (dotted)
  \foreach \n in {G, D, CI, ID}
    \draw[dotted, color=accent!60, thin] (\n.east) -- ++( 1.1cm, 0);
\end{tikzpicture}
\end{center}
\bigskip

Each arrow is non-trivial. Reading conditional independences off a graph
requires d-separation. Translating graphical structure into identification
formulas requires the back-door criterion, front-door criterion, and
do-calculus --- developed in Chapter~3. Section~5 traces Steps~1--3 in full
for the Education--Earnings DAG and previews Step~4 with the back-door
formula, whose derivation is given in Chapter~3.

% ==============================================================
\section{Summary}
% ==============================================================

\begin{enumerate}
  \item A DAG encodes the qualitative causal structure of a data-generating
        mechanism without committing to functional forms or error distributions.

  \item The three elementary path patterns --- chain, fork, and collider ---
        determine when conditioning on a middle node blocks or opens a path.
        Chains and forks are \emph{open by default} and blocked by conditioning;
        colliders are \emph{closed by default} and opened by conditioning.

  \item The d-separation criterion reads off all conditional independences
        implied by a DAG: $X$ and $Y$ are d-separated by $\mathbf{S}$
        if and only if every path between them is blocked by $\mathbf{S}$.

  \item The Markov factorization $p = \prod_i p(x_i \mid \Pa(x_i))$
        is the formal bridge between the graph and the joint distribution,
        and is the starting point for all identification derivations.

  \item Collider bias \citep{berkson1946limitations} is one of the most
        common errors in applied causal inference. Conditioning on a
        descendant of a collider opens the collider just as conditioning
        on the collider itself does.

  \item The core value of d-separation is the translation it provides: a
        qualitative causal hypothesis encoded as a DAG becomes a precise set
        of conditional independence assumptions on the joint distribution of
        the observed variables.  The IV DAG analysis in
        Section~\ref{w2:sec:ivdag} illustrates this directly --- the claim
        that $Z$ is a valid instrument is equivalent to three d-separation
        statements that can in principle be tested or exploited in data.
\end{enumerate}

% ==============================================================
\newpage
\section*{Problems}
\addcontentsline{toc}{section}{Problems}
% ==============================================================

\begin{enumerate}[label=\textbf{\arabic*.}, leftmargin=*, itemsep=10pt]

  \item \textbf{Warm-up: a single collider.}
  Consider the DAG $X \to Y \leftarrow Z$, where $X$ and $Z$ have no
  other connections.
  \begin{enumerate}[label=(\alph*)]
    \item Identify the structural role of $Y$ on the path $X \to Y \leftarrow Z$.
    \item Is $X \indep Z$?  Apply the d-separation criterion to
          the only path between $X$ and $Z$, and state which blocking rule
          applies.
    \item Is $X \indep Z \mid Y$?  Explain what happens to the
          path when $Y$ is conditioned on, and describe in one sentence the
          real-world phenomenon this illustrates.
  \end{enumerate}

  \item \textbf{d-Separation practice.}
  Consider the DAG: $A \to B \to D$, $A \to C \to D$, $B \to E$, $C \to E$.
  \begin{enumerate}[label=(\alph*)]
    \item List all paths between $A$ and $E$.  (\emph{Hint}: there are four
          paths in total; two pass through $D$.)
    \item For each path, identify the role (chain, fork, collider) of each
          intermediate node.
    \item Does $\{B, C\}$ d-separate $A$ and $E$?
    \item Does $\{D\}$ d-separate $B$ and $C$? What type of node is $D$
          on the path $B \to D \leftarrow C$?
  \end{enumerate}

  \item \textbf{Berkson's bias.}
  Suppose $X$ and $Y$ are independent standard normal variables, and let
  $S = \mathbf{1}[X + Y > 0]$ (selected into a sample).
  \begin{enumerate}[label=(\alph*)]
    \item Verify analytically that $\mathrm{Cov}(X, Y \mid S{=}1) < 0$.
    \item Draw the DAG for $(X, Y, S)$ and identify $S$ as a collider.
    \item Explain in one sentence why restricting the analysis to the
          subsample with $S=1$ biases estimates of any association between
          $X$ and $Y$, and name the type of bias this illustrates.
  \end{enumerate}

  \item \textbf{Markov factorization and collider activation.}
  Consider the DAG with edges $A \to E$, $A \to W$, $F \to E$, $E \to W$,
  where $A$ = ability, $F$ = family income, $E$ = education, $W$ = wages.
  \begin{enumerate}[label=(\alph*)]
    \item Write down the Markov factorization $p(a, f, e, w)$.
    \item Is $(F \indep W)_{\Gcal}$?
          List all paths between $F$ and $W$ and determine which are open.
    \item Is $(F \indep W \mid E)_{\Gcal}$?
          Identify the role of $E$ on each path and explain whether
          conditioning on $E$ opens or closes each one.
    \item A researcher regresses $W$ on $E$ and $F$, omitting $A$.
          Is the coefficient on $E$ a causal effect of education on wages?
          Explain using the graph.
  \end{enumerate}

  \item \textbf{Theorem~\ref{w2:thm:dsep} for the collider.}
  Consider the collider $A \to M \leftarrow B$ with Markov factorization
  $p(a, m, b) = p(a)\,p(b)\,p(m \mid a, b)$.
  \begin{enumerate}[label=(\alph*)]
    \item By marginalising over $M$, show that $A \indep B$
          in the joint distribution (i.e., $p(a,b) = p(a)\,p(b)$).
          This verifies Theorem~\ref{w2:thm:dsep} for the case
          $\mathbf{S} = \varnothing$, where
          Example~\ref{w2:ex:threesettings}(iii) established d-separation
          graphically.
    \item Show that conditioning on $M$ breaks this independence: write
          out $p(a, b \mid m)$ and explain why it does \emph{not} factorise
          into $p(a \mid m)\,p(b \mid m)$ in general.  What does this imply
          about $A$ and $B$ among hospitalised patients in the
          accident--hospitalization--cancer example?
    \item Explain in one sentence why parts~(a) and~(b) together are
          consistent with Theorem~\ref{w2:thm:dsep}.  (\emph{Hint}:
          Theorem~\ref{w2:thm:dsep} is a one-directional statement.)
  \end{enumerate}

\end{enumerate}

% ==============================================================
\newpage
\ifSubfilesClassLoaded{
  \bibliographystyle{plainnat}
  \bibliography{causal_inference}
}{}

\end{document}